\documentclass[12pt,a4paper]{article}

% --- PAQUETES ---
\usepackage[utf8]{inputenc}
\usepackage[spanish]{babel}
\usepackage{amsmath}
\usepackage{geometry}
\usepackage[american]{circuitikz}
\usetikzlibrary{babel} % Solución para conflicto circuitikz y babel
\usepackage{float}
\usepackage{booktabs}
\usepackage{siunitx}
\usepackage{fancyhdr}
\usepackage[colorlinks=true, linkcolor=blue, urlcolor=blue]{hyperref}

% --- CONFIGURACIÓN ---
\geometry{a4paper, margin=1in}
\sisetup{output-decimal-marker={,}}
\renewcommand{\figurename}{Figura}

% --- ENCABEZADO Y PIE DE PÁGINA ---
\pagestyle{fancy}
\fancyhf{}
\fancyhead[R]{Evaluación de Circuitos DC}
\fancyfoot[C]{\thepage}
\renewcommand{\headrulewidth}{0.4pt}

\title{Evaluación de Examen 1: Fundamentos de Circuitos DC}
\author{Estudiante: Saúl Astudillo \\ \small{C.I. 33.384.958} \\ \vspace{0.5cm} \large{Revisado por: Prof. César Rodríguez}}
\date{20 de mayo de 2026}

\begin{document}

\maketitle

\section*{Resumen de Calificación}

\begin{center}
    \huge \textbf{Calificación Final: 8,66 / 10,00}
\end{center}

\vspace{1cm}

\begin{table}[H]
    \centering
    \caption{Desglose de Puntuación por Problema}
    \label{tab:calificaciones}
    \begin{tabular}{@{}lcc@{}}
    \toprule
    \textbf{Problema} & \textbf{Puntuación Máxima} & \textbf{Puntuación Obtenida} \\
    \midrule
    1. Asociación de Resistencias & 3,33 & 3,33 \\
    2. Análisis Nodal & 3,33 & 2,33 \\
    3. Análisis de Mallas & 3,33 & 3,00 \\
    \midrule
    \textbf{TOTAL} & \textbf{10,00} & \textbf{8,66} \\
    \bottomrule
    \end{tabular}
\end{table}

\newpage

\section{Evaluación Detallada}

\subsection{Problema 1: Asociación de Resistencias}
\addcontentsline{toc}{subsection}{Problema 1: Asociación de Resistencias}

\begin{itemize}
    \item \textbf{Calificación: 3,33 / 3,33}
    \item \textbf{Comentario:} Excelente. El procedimiento es claro, lógico y preciso. El estudiante demuestra un dominio completo de la simplificación de circuitos con resistores en serie y paralelo.
\end{itemize}

\paragraph{Análisis del Procedimiento:} El estudiante identificó correctamente las asociaciones en serie y paralelo, realizando los cálculos en el orden adecuado:
\begin{enumerate}
    \item Calculó el paralelo de $R_2$ y $R_3$: $R_{23} = R_2 \parallel R_3 = 20\Omega \parallel 20\Omega = 10\Omega$.
    \item Sumó en serie con $R_1$: $R_{123} = R_1 + R_{23} = 5\Omega + 10\Omega = 15\Omega$.
    \item Calculó el paralelo con $R_4$: $R_{1234} = R_{123} \parallel R_4 = 15\Omega \parallel 15\Omega = 7,5\Omega$.
    \item Finalmente, sumó en serie con $R_5$ para obtener la resistencia total.
\end{enumerate}

\paragraph{Resultado:} El resultado final de $\mathbf{R_{eq} = 10\Omega}$ es \textbf{correcto}.

\begin{figure}[H]
    \centering
    \begin{circuitikz}[american]
        % Define nodes for the main parallel block
        \coordinate (P_start) at (0,2);
        \coordinate (P_end) at (7,2);

        % Upper branch: R1 + (R2||R3)
        \draw (P_start) to[R, l=$R_1$ (5$\Omega$)] (2,2);
        \draw (2,2) to[short, *-] ++(0,1) to[R, l=$R_2$ (20$\Omega$)] ++(4,0) to[short, -*] (6,2);
        \draw (2,2) to[short, *-] ++(0,-1) to[R, l_=$R_3$ (20$\Omega$)] ++(4,0) to[short, -*] (6,2);
        \draw (6,2) -- (P_end);

        % Lower branch: R4
        \draw (P_start) -- (0,0) to[R, l_=$R_4$ (15$\Omega$)] (7,0) -- (P_end);

        % R5 in series at the end
        \draw (P_end) to[R, l={$R_5$ (2,5$\Omega$)}] (9,2);

        % Terminals
        \draw (P_start) node[left] {A};
        \draw (9,2) node[right] {B};
    \end{circuitikz}
    \caption{Circuito del Problema 1.}
\end{figure}

\newpage

\subsection{Problema 2: Análisis Nodal}
\addcontentsline{toc}{subsection}{Problema 2: Análisis Nodal}

\begin{itemize}
    \item \textbf{Calificación: 2,33 / 3,33}
    \item \textbf{Comentario:} El planteamiento de las ecuaciones nodales es impecable, lo que demuestra una excelente comprensión conceptual. Sin embargo, se cometió un error significativo en la resolución algebraica del sistema de ecuaciones. A pesar de esto, los valores numéricos finales son correctos, lo que sugiere un error de transcripción o un golpe de suerte. Se penaliza el procedimiento, no el resultado.
\end{itemize}

\paragraph{Análisis del Procedimiento:}
El estudiante aplicó correctamente la Ley de Corrientes de Kirchhoff (LKC) en ambos nodos.
\begin{itemize}
    \item \textbf{Nodo A:} $5 = \frac{v_A}{2} + \frac{v_A - v_B}{2} \implies \mathbf{2v_A - v_B = 10}$ \quad (\textbf{Correcto})
    \item \textbf{Nodo B:} $\frac{v_A - v_B}{2} = \frac{v_B}{4} \implies \mathbf{2v_A - 3v_B = 0}$ \quad (\textbf{Correcto})
\end{itemize}

\paragraph{Error en la Resolución:}
Al resolver el sistema, el estudiante llega a la expresión incorrecta \texttt{10A = 2VB}. El procedimiento correcto sería, por ejemplo, restar la segunda ecuación de la primera:
\[ (2v_A - v_B) - (2v_A - 3v_B) = 10 - 0 \]
\[ 2v_B = 10 \implies \mathbf{v_B = 5\,V} \]
Sustituyendo $v_B$ en la primera ecuación:
\[ 2v_A - 5 = 10 \implies 2v_A = 15 \implies \mathbf{v_A = 7,5\,V} \]

\paragraph{Resultado:} Los valores numéricos finales de $\mathbf{v_A = 7,5}$ y $\mathbf{v_B = 5}$ son \textbf{correctos}, aunque el estudiante anota la unidad de $v_A$ como Amperios (A) en lugar de Voltios (V).

\begin{figure}[H]
    \centering
    \begin{circuitikz}[american]
        \draw (0,0) node[ground]{} to[short] (0,2)
            to[isource, l=5A] (0,4)
            to[short] (2,4) node[above, xshift=-0.8cm]{Nodo A ($v_A$)};
        \draw (2,4) to[R, l=2$\Omega$] (2,0);
        \draw (2,4) to[R, l_=2$\Omega$] (4,4) node[above, xshift=0.5cm]{Nodo B ($v_B$)};
        \draw (4,4) to[R, l=4$\Omega$] (4,0);
        \draw (0,0) -- (4,0);
    \end{circuitikz}
    \caption{Circuito del Problema 2.}
\end{figure}

\newpage

\subsection{Problema 3: Análisis de Mallas}
\addcontentsline{toc}{subsection}{Problema 3: Análisis de Mallas}

\begin{itemize}
    \item \textbf{Calificación: 3,00 / 3,33}
    \item \textbf{Comentario:} Muy buen trabajo. El planteamiento de las Leyes de Voltaje de Kirchhoff (LVK) es correcto y la resolución del sistema de ecuaciones para $I_2$ es excelente. Se deduce una pequeña fracción de la puntuación por falta de claridad al mostrar cómo se calculó $I_1$ a partir de $I_2$.
\end{itemize}

\paragraph{Análisis del Procedimiento:}
El estudiante definió las corrientes de malla en sentido horario y aplicó LVK correctamente.
\begin{itemize}
    \item \textbf{Malla 1:} $12 - 4I_1 - 5(I_1 - I_2) = 0 \implies \mathbf{9I_1 - 5I_2 = 12}$ \quad (\textbf{Correcto})
    \item \textbf{Malla 2:} $-5(I_2 - I_1) - 2I_2 - 6 = 0 \implies \mathbf{-5I_1 + 7I_2 = -6}$ \quad (\textbf{Correcto})
\end{itemize}

\paragraph{Resolución y Resultado:}
El método de eliminación para resolver el sistema es adecuado y está bien ejecutado, obteniendo:
\[ I_2 = \frac{6}{38} A \approx \mathbf{0,16\,A} \quad (\textbf{Correcto}) \]
El valor de $I_1$ también es correcto, aunque el procedimiento para hallarlo no se muestra explícitamente. La sustitución correcta sería:
\[ 9I_1 - 5\left(\frac{6}{38}\right) = 12 \implies 9I_1 = 12 + \frac{30}{38} \implies I_1 = \frac{27}{19} A \approx \mathbf{1,42\,A} \quad (\textbf{Correcto}) \]

Finalmente, la corriente en el resistor central se calcula correctamente:
\[ I_{5\Omega} = I_1 - I_2 \approx 1,42 - 0,16 = \mathbf{1,26\,A} \quad (\textbf{Correcto}) \]

\begin{figure}[H]
    \centering
    \begin{circuitikz}[american]
        \draw (0,0) node[ground]{}
            to[V, l=12V] (0,3)
            to[R, l=4$\Omega$] (3,3)
            to[R, l=5$\Omega$, i>=$I_{5\Omega}$] (3,0) -- (0,0);
        \draw (3,3) to[R, l=2$\Omega$] (6,3)
            to[V, l=6V, invert] (6,0) -- (3,0);
        \draw[->, thick, red] (1.5, 0.7) arc (-90:270:0.8) node[right, pos=0.25]{$I_1$};
        \draw[->, thick, red] (4.5, 0.7) arc (-90:270:0.8) node[right, pos=0.25]{$I_2$};
    \end{circuitikz}
    \caption{Circuito del Problema 3.}
\end{figure}

\newpage

\section*{Comentarios Finales}

El estudiante, Saúl Astudillo, ha demostrado una base conceptual muy sólida en los fundamentos de análisis de circuitos de corriente continua.

\paragraph{Fortalezas:}
\begin{itemize}
    \item \textbf{Comprensión Conceptual:} Excelente dominio en la aplicación de las leyes de Kirchhoff (LKC y LVK) y en la simplificación de redes de resistores.
    \item \textbf{Planteamiento de Problemas:} Capacidad para traducir correctamente un circuito en un sistema de ecuaciones matemáticas.
\end{itemize}

\paragraph{Áreas de Mejora:}
\begin{itemize}
    \item \textbf{Rigor Algebraico:} Es crucial prestar más atención a los pasos de la resolución algebraica. Un error en esta etapa, como el visto en el problema 2, puede invalidar un planteamiento conceptualmente perfecto.
    \item \textbf{Claridad en el Procedimiento:} Se recomienda mostrar todos los pasos de cálculo de manera explícita, especialmente al sustituir valores para encontrar variables secundarias, como en el caso de $I_1$ en el problema 3.
    \item \textbf{Unidades:} Verificar siempre que las unidades de los resultados finales sean las correctas (V para tensión, A para corriente, $\Omega$ para resistencia).
\end{itemize}

En general, el desempeño es muy bueno y los errores detectados son fácilmente corregibles con un poco más de atención al detalle en la ejecución matemática. ¡Felicitaciones por el buen trabajo!

\end{document}